% Nouveau EGA V (cote 25): archive cote 25, site page 31 (PDF page 32).
% typescript. Draft compared: none.
% AI-transcribed by Claude Opus 5.5 (claude-opus-5-5, adaptive thinking, effort high), Message Batches API, batch ega25b (msgbatch_01A7n5YuXubUPNK1zsepi8RG), 2026-11-13; not yet checked by eye.
\scanpage{25/31}
\mspage{\unsure{27}}

8. \uline{Dimension de l'ensemble des hyperplans exceptionnels}.

8.1.Dans les n°s précédents, et notamment 2 et 3, nous avons donné des
énoncés affirmant que l'ensemble des \hand{$\xi\in$} $\overset{\hand{$\vee$}}{P}$ tels que Y ait une certaine
propriété P est constructible, et contient le point générique\inserted{\unsure{$\overset{\vee}{\ }$}}\hand{/}, \xout{xxxx} ou
encore que l'ensemble $Z_P$ des \hand{$\xi\in$} $\overset{\hand{$\vee$}}{P}$ ``exceptionnels pour P'' est construc-%
tible, et est rare, i.e. son adhérence est de\inserted{co}dimension $\unsure{\geq}$ \hand{1}\deleted{r (où r=dim $\check{P}$}) \margin{\unsure{.}}
(NB on suppose S=Spec(k)). Dans certains cas on peut préciser cet énoncé
en donnant une meilleure m\over{aj}{\unsure{in}}oration de cette\inserted{co}dimension, ce qui est impor-%
tant \xout{pour}\inserted[typed]{dans} diverses questions. Par exemple, lorsqu'on sait que cette\inserted{co}dimen
sion est $\geq$ \deleted{\unsure{r}}-2, il s'ensuit qu'une droite ``assez générale'' D de $\overset{\hand{$\vee$}}{P}$
ne rencontre pas $Z_P$ , d'où l'existence (si k est infini) de ``pinceaux
linéaires'' de sections hyperplanes \xout{(Xx)xxxxx}\hand{\unsure{Dt}} $Y_{\hand{\xi}}$ (\hand{$\xi$} un point géométri
que de D) qui sont toutes la propriété P (cf n° \deleted{suivant}\inserted{\illegible{1}} pour des exem-%
ples. - Du point de vue rédaction, comme les résultats du présent n°
précisent certains résultats des n°s antérieurs, la question se pose s'il
y a bien lieu de faire ce ratrappage dans un n° à part, ou de donner la
version précisée à mesure. Redactor decidetur.

\deleted{8.2. Soit Z une partie \deleted{constructible}\inserted{telle que} de X, et considérons l'ensemble \margin{Z}
des \quad P tels que son image inverse Z dans Y soit de dimension
\inserted[typed]{strictement}plus grande (i.e. son adhérence \inserted[typed]{dans Y}de dimension \inserted[typed]{str.}plus grande) que celle de
la fibre générique (i.e. soit égale à dim Z si dim $\overline{f(Z)}$ \quad 0, resp.
soit non vide si dim $\overline{f(Z)}$ \hand{\unsure{(}} 0 ). Alors Z est de codim\xout{\unsure{mm}}nsion \quad 2 i.e.
de dimension \quad r-2. (Cela résulte aussitôt de la théorie de la dimension
d'un morphisme appliqué à $\underline{Y}$ \quad P ). Bien entendu, cet énoncé peut se
\inserted[typed]{en termes de supp F}moduler (et est équivalent à l'énoncé modulé), et l'appliquant aux com-%
posantes irréductibles géométriques de T, on trouve aussi que l'ensemble
Z des \quad tels \xout{xxx} qu'il existe une composante irr de}

\mspage{31}
